\documentclass[11pt]{article}

% --- UNIVERSAL PREAMBLE BLOCK ---
\usepackage[a4paper, left=1.06cm, top=1.7cm, right=1.06cm, bottom=1.5cm]{geometry}
\usepackage{fontspec}

\usepackage[bidi=basic, provide=*]{babel}

\babelprovide[main, import, onchar=ids fonts]{spanish}
\babelprovide[import, onchar=ids fonts]{english}

% Set default/Latin font to Sans Serif in the main (rm) slot
\babelfont{rm}{Noto Sans}

% Add because we replaced the legacy list behavior
\usepackage{enumitem}
\setlist[itemize]{label=-}
% --- END UNIVERSAL PREAMBLE BLOCK ---

\usepackage{hyperref}
\setlength{\parindent}{0pt}

\begin{document}
\begin{center}
    \textbf{\LARGE Miguel Angel Bravo Vidales}\\ 
    \vspace{4pt}
    \hrulefill
\end{center}

\begin{center}
    San Luis Potosí, SLP, México \textbullet \ x@mackaber.me \textbullet \ +52 444 444 7397 \textbullet \ linkedin.com/in/mackaber \textbullet \ github.com/mackaber
\end{center}

\vspace{0.5pt}

\begin{center}
    \textbf{Educación}
\end{center}
\textbf{Tecnológico de Monterrey (ITESM)} \hfill Monterrey, NL

Maestría en Ciencias de la Computación (Enfoque en Inteligencia Artificial y Aprendizaje Automático) \hfill Dic 2020

\vspace{10pt}

\textbf{Tecnológico de Monterrey (ITESM)} \hfill San Luis Potosí, SLP

Ingeniería en Tecnologías de Información y Comunicaciones \hfill Dic 2013

\vspace{12pt}

\begin{center}
    \textbf{Experiencia}
\end{center}

\textbf{Tecnológico de Monterrey (ITESM)} \hfill San Luis Potosí, SLP

\textbf{Profesor de Preparatoria} \hfill Agosto 2025 -- Presente
\begin{itemize}[noitemsep, topsep=0pt, partopsep=0pt, parsep=0pt]
    \item Impartí clases de desarrollo de software a estudiantes de preparatoria, cubriendo fundamentos de programación, construcción de lógica y diseño práctico de aplicaciones.
\end{itemize}

\vspace{8pt}

\textbf{Encora, Inc.} \hfill Remoto

\textbf{Instructor de AWS re/Start} \hfill Febrero 2025 -- Junio 2025
\begin{itemize}[noitemsep, topsep=0pt, partopsep=0pt, parsep=0pt]
    \item Impartí capacitación integral en cómputo en la nube a un grupo de estudiantes, cubriendo servicios principales de AWS, redes, seguridad y fundamentos de Linux/Python.
    \item Guié a los estudiantes en laboratorios prácticos para prepararlos para la certificación AWS Certified Cloud Practitioner y su transición a roles de entrada en cómputo en la nube.
\end{itemize}

\vspace{8pt}

\textbf{Spark Tech Staff} \hfill Enero 2025 -- Presente
\begin{itemize}[noitemsep, topsep=0pt, partopsep=0pt, parsep=0pt]
    \item Mentoricé a nuevas contrataciones en estructuras de datos y algoritmos mediante el programa interno Algorithm School Program, acelerando su preparación para asignación con clientes.
\end{itemize}

\vspace{8pt}

\textbf{Desarrollador Full Stack} \hfill Enero 2021 -- 2024
\begin{itemize}[noitemsep, topsep=0pt, partopsep=0pt, parsep=0pt]
    \item Construí herramientas de desarrollo impulsadas por IA usando Python, LangChain, Streamlit y la API de OpenAI, automatizando la generación de aplicaciones plantilla en React y React Native (Expo).
    \item Mantuve y optimicé una aplicación heredada de lealtad para viajes (Arrivia), implementando monitoreo de recursos con Azure AppInsights para identificar y resolver tareas de alto consumo.
    \item Desarrollé pipelines robustos de procesamiento de datos para una plataforma de telesalud (MDLive), procesando expedientes de pacientes y construyendo sistemas de incentivos para médicos con Ruby on Rails, Redis y RabbitMQ.
\end{itemize}

\vspace{8pt}

\textbf{DEV.f} \hfill Remoto

\textbf{Profesor de Desarrollo Web} \hfill Agosto 2021 -- Marzo 2024
\begin{itemize}[noitemsep, topsep=0pt, partopsep=0pt, parsep=0pt]
    \item Impartí bootcamps de desarrollo web full-stack de 6 meses cubriendo HTML, CSS, JavaScript, Node.js y React.
    \item Guié a estudiantes sin experiencia previa en desarrollo de software mediante proyectos prácticos para construir portafolios listos para la industria.
\end{itemize}

\vspace{8pt}

\textbf{Continental} \hfill San Luis Potosí, SLP

\textbf{Desarrollador Full Stack (Freelancer)} \hfill Enero 2014 -- Diciembre 2017
\begin{itemize}[noitemsep, topsep=0pt, partopsep=0pt, parsep=0pt]
    \item Desarrollé y di mantenimiento a la aplicación Android para una plataforma de administración de transportadores.
    \item Construí un servicio de integración de escritorio en Ruby para conectar nuevas interfaces con aplicaciones legadas existentes.
\end{itemize}

\vspace{8pt}

\textbf{Winero} \hfill San Luis Potosí, SLP

\textbf{Cofundador, CTO \& Desarrollador Full Stack} \hfill Julio 2013 -- Diciembre 2016
\begin{itemize}[noitemsep, topsep=0pt, partopsep=0pt, parsep=0pt]
    \item Cofundé una plataforma de programas de lealtad, diseñando la arquitectura de datos, el sistema backend (Ruby on Rails, PostgreSQL) y el panel de analítica de rendimiento.
\end{itemize}

\vspace{12pt}

\begin{center}
    \textbf{Liderazgo y Actividades}
\end{center}

\textbf{Comunidades y Eventos Tech} \hfill México

\textbf{Speaker, Mentor y Cofundador} \hfill 2013 -- Presente
\begin{itemize}[noitemsep, topsep=0pt, partopsep=0pt, parsep=0pt]
    \item Participé en más de 20 hackathons (ganando 2 y mentoreando 10), incluyendo Startup Weekend, AngelHack y Talent Hackathon.
    \item Impartí más de 15 charlas y talleres sobre ingeniería de software e IA en eventos tecnológicos como TalentLand, Nerdearla y NerdyTalks.
    \item Cofundé ``Comunidad Código,'' una asociación local de desarrolladores, y contribuí activamente a la iniciativa Tuna Valley.
\end{itemize}

\vspace{12pt}

\begin{center}
    \textbf{Habilidades, Certificaciones e Intereses}
\end{center}

\textbf{Certificaciones:} AWS Solutions Architect (2025), Azure AI Engineer Associate (AI-102) (2024, renovada 2025)

\textbf{Habilidades técnicas:} Ruby, Python, JavaScript, Java, SQL, Ruby on Rails, Node.js, React, TensorFlow, Docker, AWS, Azure, Git, LangChain, Redis, RabbitMQ

\textbf{Idiomas:} Español (Nativo), Inglés (Fluido)

\end{document}